\subsection{Structural Extraction and Evaluation}

A Python string is a sequence. That means it supports the same basic positional operations introduced earlier: indexing, slicing, length checking, membership testing, and iteration. The difference is that every component extracted from a string is itself a string.

Python does not expose a separate character type. A single character is represented as a \texttt{str} object of length \texttt{1}.

\begin{verbatim}
text = "Ni!"

print(f"text: {text!r}")
print(f"type(text): {type(text)}")
print(f"len(text): {len(text)}")

print()
for ch in text:
    print(f"ch={ch!r:<4} type={type(ch)} len={len(ch)}")
\end{verbatim}

Possible output:

\begin{verbatim}
text: 'Ni!'
type(text): <class 'str'>
len(text): 3

ch='N'  type=<class 'str'> len=1
ch='i'  type=<class 'str'> len=1
ch='!'  type=<class 'str'> len=1
\end{verbatim}

Some string operations used in this section are visible through \texttt{dir()}.

\begin{verbatim}
text = "Nobody expects the Spanish Inquisition"

selected_names = [
    "__getitem__",
    "__len__",
    "__contains__",
    "format",
    "format_map",
    "center",
    "ljust",
    "rjust",
]

for name in selected_names:
    print(f"{name:<14} {name in dir(text)}")
\end{verbatim}

Possible output:

\begin{verbatim}
__getitem__    True
__len__        True
__contains__   True
format         True
format_map     True
center         True
ljust          True
rjust          True
\end{verbatim}

The syntax \texttt{text[index]} and \texttt{text[start:stop:step]} is ordinary Python syntax, but internally it is still item access on the string object.

\subsubsection{Indexing and Slicing: Characters as One-Character String Objects}

Indexing extracts one component from a string.

\begin{verbatim}
text = "D'oh!"

print(f"text: {text!r}")
print(f"text[0]: {text[0]!r}")
print(f"text[1]: {text[1]!r}")
print(f"text[2]: {text[2]!r}")
print(f"text[3]: {text[3]!r}")
print(f"text[4]: {text[4]!r}")
\end{verbatim}

Possible output:

\begin{verbatim}
text: "D'oh!"
text[0]: 'D'
text[1]: "'"
text[2]: 'o'
text[3]: 'h'
text[4]: '!'
\end{verbatim}

Each extracted component is a string object.

\begin{verbatim}
text = "D'oh!"

ch = text[0]

print(f"ch: {ch!r}")
print(f"type(ch): {type(ch)}")
print(f"len(ch): {len(ch)}")
\end{verbatim}

Possible output:

\begin{verbatim}
ch: 'D'
type(ch): <class 'str'>
len(ch): 1
\end{verbatim}

Negative indexing counts from the right.

\begin{verbatim}
text = "D'oh!"

print(f"text[-1]: {text[-1]!r}")
print(f"text[-2]: {text[-2]!r}")
print(f"text[-3]: {text[-3]!r}")
print(f"text[-4]: {text[-4]!r}")
print(f"text[-5]: {text[-5]!r}")
\end{verbatim}

Possible output:

\begin{verbatim}
text[-1]: '!'
text[-2]: 'h'
text[-3]: 'o'
text[-4]: "'"
text[-5]: 'D'
\end{verbatim}

An invalid index raises \texttt{IndexError}.

\begin{verbatim}
text = "D'oh!"

try:
    print(text[42])
except IndexError as err:
    print("indexing failed")
    print(f"message: {err}")
\end{verbatim}

Possible output:

\begin{verbatim}
indexing failed
message: string index out of range
\end{verbatim}

Slicing extracts a substring.

\begin{verbatim}
text = "Nobody expects the Spanish Inquisition"

print(f"text[0:6]:   {text[0:6]!r}")
print(f"text[7:14]:  {text[7:14]!r}")
print(f"text[19:26]: {text[19:26]!r}")
print(f"text[27:]:   {text[27:]!r}")
\end{verbatim}

Possible output:

\begin{verbatim}
text[0:6]:   'Nobody'
text[7:14]:  'expects'
text[19:26]: 'Spanish'
text[27:]:   'Inquisition'
\end{verbatim}

The result of a string slice is also a \texttt{str}.

\begin{verbatim}
text = "No soup for you!"

part = text[3:7]

print(f"part: {part!r}")
print(f"type(part): {type(part)}")
print(f"len(part): {len(part)}")
\end{verbatim}

Possible output:

\begin{verbatim}
part: 'soup'
type(part): <class 'str'>
len(part): 4
\end{verbatim}

An empty slice is still a string.

\begin{verbatim}
text = "abc"

part = text[10:20]

print(f"part: {part!r}")
print(f"type(part): {type(part)}")
print(f"len(part): {len(part)}")
\end{verbatim}

Possible output:

\begin{verbatim}
part: ''
type(part): <class 'str'>
len(part): 0
\end{verbatim}

So the structural rule is consistent:

\begin{quote}
Indexing a string returns a one-character \texttt{str}. Slicing a string returns a possibly empty \texttt{str}.
\end{quote}

\subsubsection{Substring Evaluation Mechanisms: Step-Based Slice Offsets (\texttt{[start:stop:step]}) vs. Manual Pointer Offsets}

A slice uses three values:

\begin{itemize}
    \item \texttt{start}: where extraction begins;
    \item \texttt{stop}: where extraction stops, not included;
    \item \texttt{step}: how far Python moves after taking a component.
\end{itemize}

\begin{verbatim}
text = "abcdefghij"

print(f"text[0:10:1]: {text[0:10:1]!r}")
print(f"text[0:10:2]: {text[0:10:2]!r}")
print(f"text[1:10:2]: {text[1:10:2]!r}")
print(f"text[2:9:3]:  {text[2:9:3]!r}")
\end{verbatim}

Possible output:

\begin{verbatim}
text[0:10:1]: 'abcdefghij'
text[0:10:2]: 'acegi'
text[1:10:2]: 'bdfhj'
text[2:9:3]:  'cfi'
\end{verbatim}

The slice \texttt{text[2:9:3]} means:

\begin{enumerate}
    \item start at index \texttt{2};
    \item stop before index \texttt{9};
    \item take every third component.
\end{enumerate}

This can be shown manually.

\begin{verbatim}
text = "abcdefghij"

print("manual positions for text[2:9:3]")
print()

for index in range(2, 9, 3):
    ch = text[index]
    print(f"index={index} ch={ch!r}")

print()
print(f"slice result: {text[2:9:3]!r}")
\end{verbatim}

Possible output:

\begin{verbatim}
manual positions for text[2:9:3]

index=2 ch='c'
index=5 ch='f'
index=8 ch='i'

slice result: 'cfi'
\end{verbatim}

This is not C pointer arithmetic. In C, a programmer may think in terms of a starting address plus offsets into memory. In Python, the programmer describes positions in a sequence. Python evaluates the requested positions and constructs the resulting string value.

A manual reconstruction makes the difference visible.

\begin{verbatim}
text = "abcdefghij"

manual = ""

for index in range(2, 9, 3):
    manual = manual + text[index]

sliced = text[2:9:3]

print(f"manual: {manual!r}")
print(f"sliced: {sliced!r}")
print(f"same:   {manual == sliced}")
\end{verbatim}

Possible output:

\begin{verbatim}
manual: 'cfi'
sliced: 'cfi'
same:   True
\end{verbatim}

The two results are equal, but the slice is the direct Python expression for the operation. The manual loop is useful for understanding, not for ordinary substring extraction.

A negative step walks backward.

\begin{verbatim}
text = "abcdef"

print(f"text[::-1]:    {text[::-1]!r}")
print(f"text[4:1:-1]:  {text[4:1:-1]!r}")
print(f"text[5:0:-2]:  {text[5:0:-2]!r}")
\end{verbatim}

Possible output:

\begin{verbatim}
text[::-1]:    'fedcba'
text[4:1:-1]:  'edc'
text[5:0:-2]:  'fdb'
\end{verbatim}

The same operation can again be expanded into indexes.

\begin{verbatim}
text = "abcdef"

print("manual positions for text[5:0:-2]")
print()

for index in range(5, 0, -2):
    print(f"index={index} ch={text[index]!r}")

print()
print(f"slice result: {text[5:0:-2]!r}")
\end{verbatim}

Possible output:

\begin{verbatim}
manual positions for text[5:0:-2]

index=5 ch='f'
index=3 ch='d'
index=1 ch='b'

slice result: 'fdb'
\end{verbatim}

The \texttt{step} cannot be zero.

\begin{verbatim}
text = "abc"

try:
    part = text[0:3:0]
    print(part)
except ValueError as err:
    print("slicing failed")
    print(f"message: {err}")
\end{verbatim}

Possible output:

\begin{verbatim}
slicing failed
message: slice step cannot be zero
\end{verbatim}

A zero step would mean ``move by no positions'', so Python refuses the slice.

The practical rule is:

\begin{quote}
A string slice describes a positional extraction pattern. It is not manual memory traversal. Python chooses the indexes, reads the corresponding string components, and returns a string result.
\end{quote}

\subsubsection{String Formatting Paradigms: Variable Injection via Modern f-Strings}

String formatting builds a string by combining literal text with values computed by the program. Modern Python code commonly uses f-strings.

An f-string begins with \texttt{f} before the opening quote. Expressions inside braces are evaluated and inserted into the resulting string.

\begin{verbatim}
name = "Kramer"
score = 42

line = f"{name} has score {score}."

print(line)
\end{verbatim}

Possible output:

\begin{verbatim}
Kramer has score 42.
\end{verbatim}

The expression inside braces can be a variable, an arithmetic expression, or an indexing operation.

\begin{verbatim}
name = "Homer"
numbers = [10, 20, 12]

line1 = f"name: {name}"
line2 = f"first letter: {name[0]}"
line3 = f"total: {numbers[0] + numbers[1] + numbers[2]}"

print(line1)
print(line2)
print(line3)
\end{verbatim}

Possible output:

\begin{verbatim}
name: Homer
first letter: H
total: 42
\end{verbatim}

The inserted value is converted to text. For debugging, \texttt{!r} asks for the representation form, which is often more explicit than the printed form.

\begin{verbatim}
text = "C:\new\test"
raw_text = r"C:\new\test"

print(f"text printed:    {text}")
print(f"text repr:       {text!r}")

print()
print(f"raw printed:     {raw_text}")
print(f"raw repr:        {raw_text!r}")
\end{verbatim}

Possible output:

\begin{verbatim}
text printed:    C:
ew	est
text repr:       'C:\new\test'

raw printed:     C:\new\test
raw repr:        'C:\\new\\test'
\end{verbatim}

The \texttt{!r} form is useful when teaching strings because it exposes quotes, backslashes, newline escapes, and tabs more clearly.

F-strings also support alignment.

\begin{verbatim}
names = ["Jerry", "Elaine", "George", "Kramer"]

for index, name in enumerate(names):
    print(f"{index:>2} | {name:<10} | len={len(name):>2}")
\end{verbatim}

Possible output:

\begin{verbatim}
 0 | Jerry      | len= 5
 1 | Elaine     | len= 6
 2 | George     | len= 6
 3 | Kramer     | len= 6
\end{verbatim}

The formatting fields above mean:

\begin{itemize}
    \item \texttt{:>2}: align to the right in a width of \texttt{2};
    \item \texttt{:<10}: align to the left in a width of \texttt{10};
    \item \texttt{:>2}: align to the right in a width of \texttt{2}.
\end{itemize}

Floating-point values can be formatted with explicit precision.

\begin{verbatim}
value = 10 / 3

print(f"default: {value}")
print(f"two decimals: {value:.2f}")
print(f"four decimals: {value:.4f}")
\end{verbatim}

Possible output:

\begin{verbatim}
default: 3.3333333333333335
two decimals: 3.33
four decimals: 3.3333
\end{verbatim}

The older \texttt{str.format()} method is also still common.

\begin{verbatim}
name = "Elaine"
score = 42

line = "{} has score {}.".format(name, score)

print(line)
\end{verbatim}

Possible output:

\begin{verbatim}
Elaine has score 42.
\end{verbatim}

It also supports named fields.

\begin{verbatim}
line = "{name:<10} | score={score:>3}".format(
    name="George",
    score=42,
)

print(line)
\end{verbatim}

Possible output:

\begin{verbatim}
George     | score= 42
\end{verbatim}

For most ordinary code, f-strings are clearer because the expression appears directly where its value will be inserted.

\begin{verbatim}
name = "Bart"
age = 10

print(f"{name} is {age} years old.")
print("{} is {} years old.".format(name, age))
\end{verbatim}

Possible output:

\begin{verbatim}
Bart is 10 years old.
Bart is 10 years old.
\end{verbatim}

The practical rule is:

\begin{quote}
Use f-strings when building readable output from current variables and expressions. Use explicit alignment, precision, and \texttt{!r} when the exact textual representation matters.
\end{quote}